\section{Jagtetiske regler for kronvildt}

\begin{frame}
  \vspace{3em}
  \begin{center}
    {\LARGE \textbf{Jagtetiske regler for kronvildt}
    }
    \\
    \vspace{1em}
    Vildtforvaltningsrådet \\
    https://sgavmst.dk/media/vcupseet/de-jagtetiske-regler-for-kronvildt.pdf
  \end{center}
\end{frame}

\begin{frame}{Formål \& Ansvar}
\begin{itemize}
\item Sikre bred offentlig accept af kronvildtjagt – grundlag for fortsat jagt uden stramme restriktioner
\item Regler udarbejdet af Vildtforvaltningsrådet
\item Den enkelte jægers ansvar: Værne om kronvildt og jagten ved at følge reglerne
\item Fokus: Naturlig bestand, nabohensyn, bæredygtighed, ren skydning og minimal forstyrrelse
\end{itemize}
\end{frame}

\begin{frame}{Afskydningspolitik}
\begin{itemize}
\item Tilrettelæg afskydning for naturlig køns- og aldersfordeling
\item Størst afskydning blandt unge dyr (kalve + unghjorte/unghinder)
\item Jævn fordeling mellem handyr og hundyr
\item Indskrænk septemberjagten på hjorte mest muligt → giver ro i brunsten
\item Undgå ubalance (mange hjorte – få hinder/kalve)
\end{itemize}
\end{frame}

\begin{frame}{Nabohensyn \& Terrænstørrelse}
\begin{itemize}
\item Afskydning skal stå i rimeligt forhold til terrænets størrelse
\item Sigt mod at vildt falder på eget område (kronvildt går ofte 100 m+)
\item Ikke flere end 1 stk. pr. sæson på arealer mindre end 25 ha
\item På små/smalle arealer: Indgå naboaftaler eller fælles jagt/forståelse
\item Skel mellem kerneområder (sædvanlig ophold) og vandringsområder
\end{itemize}
\end{frame}

\begin{frame}{Bæredygtig jagtudnyttelse}
\begin{itemize}
\item Bæredygtighed vurderes lokalt – ikke nationalt
\item Typisk max 1 stk. pr. 25–50 ha pr. sæson (lavere på større ejendomme)
\item Kerneområder (ro, dækning, føde) tåler højere afskydning pr. ha
\item Yderområder: Lavere afskydning for at etablere fast bestand
\item Kronvildt har store revirer → elastisk bæredygtighedsbegreb
\end{itemize}
\end{frame}

\begin{frame}{Skydning \& Skydefærdigheder}
\begin{itemize}
\item Vær sikker på, hvad du skyder på – hold hovedet koldt ved dyr i rudel
\item Foretræk sideskud – skyd helst på dyr i ro
\item Ingen chancebetonede skud – træn på skive + bevægeligt mål
\item Kend dine begrænsninger – træn også luftbøsse/bevægelige mål/tørtræning
\item Brug hochsitz/hochstand når muligt
\end{itemize}
\end{frame}

\begin{frame}{Anskydninger \& Eftersøgning}
\begin{itemize}
\item Vær meget opmærksom på skudtegn
\item Skyd ikke til friskt dyr, hvis tidligere skudt dyr ikke ligger synligt forendt
\item Brug altid schweisshund ved tvivl/anskydning
\item Afmærk præcist ved tvivl – tilkald erfaren schweisshundefører
\item Hjælp hundefører med gode observationer – undgå flere skud i samme område
\item Mål: Minimal antal anskydninger – skylder vi kronvildtet og hinanden
\end{itemize}
\end{frame}

\begin{frame}{Jagtudøvelse \& Forstyrrelser}
\begin{itemize}
\item Ingen tryk-/drivjagter i september (brunsten)
\item Skyd ikke diegivende hinder i oktober (kalve dier stadig)
\item Min. 3 uger mellem tryk-/drivjagter
\item Max 3–4 fællesjagter pr. sæson på samme revir
\item Begræns hurtigtgående store hunde til drivjagt
\item Pürsch foretrækkes – undgå simultan pürsch over hele terræn
\item Frihold områder under fællesjagter – kronvildt kræver ro (især i etableringsfase)
\end{itemize}
\end{frame}

\begin{frame}{Dialog \& Regionale Krondyrgrupper}
\begin{itemize}
\item Planlæg afskydning lokalt sammen med naboer/lodsejere
\item Indgå i (formelle/uformelle) samarbejder om bestand og skader
\item Bidrag til fælles ansvarsfølelse
\item Regionale Krondyrgrupper understøtter efterlevelse af reglerne
\item Bred repræsentation → bedre dialog og forvaltning
\end{itemize}

\bigskip
\textbf{Sammenfatning: Kronvildtjagt kræver særligt hensyn – forvaltes med respekt, samarbejde og selvbegrænsning.}
\end{frame}